\chapter{Introduction, Literature Review and Problem Formulation}
\label{chap:introduction}

The exponential growth of wireless communication devices and the insatiable demand for higher data rates have created unprecedented challenges for modern telecommunication networks. With projections indicating over 50 billion Internet of Things (IoT) devices and mobile users by 2030, coupled with the emergence of bandwidth-intensive applications such as ultra-high-definition video streaming, augmented reality, autonomous vehicles, and industrial automation, wireless network infrastructure faces enormous pressure. Fifth Generation (5G) networks and the upcoming Sixth Generation (6G) systems must address these challenges while ensuring spectral efficiency, energy efficiency, low latency, and fairness across heterogeneous user populations.

Traditional Orthogonal Multiple Access (OMA) schemes, including Frequency Division Multiple Access (FDMA), Time Division Multiple Access (TDMA), Code Division Multiple Access (CDMA), and Orthogonal Frequency Division Multiple Access (OFDMA), allocate exclusive time-frequency resources to each user. While this orthogonal allocation eliminates intra-cell interference, it fundamentally limits spectral efficiency as the number of users grows. Cell-edge users experiencing poor channel conditions consume disproportionate resources to meet minimum Quality of Service (QoS) requirements, further constraining system capacity. This limitation has motivated the research community to explore non-orthogonal transmission strategies that can accommodate multiple users on the same time-frequency resources.

\section{Non-Orthogonal Multiple Access: A Paradigm Shift}

Power-Domain Non-Orthogonal Multiple Access (PD-NOMA) has emerged as a transformative technology to overcome the fundamental limitations of orthogonal access schemes. Unlike OMA, PD-NOMA allows multiple users to share identical time-frequency resources by exploiting the power domain through an ingenious combination of three key techniques. First, asymmetric power allocation assigns higher transmit power to users with weaker channel conditions (cell-edge users) while allocating lower power to users with stronger channels (cell-center users). Second, superposition coding enables the base station to transmit a composite signal containing information for multiple users simultaneously on the same Physical Resource Block (PRB). Third, Successive Interference Cancellation (SIC) allows stronger users to decode and remove signals intended for weaker users before decoding their own, thereby canceling interference.

This approach achieves significant spectral efficiency gains, typically demonstrating 30-50\% improvement over OFDMA in dense deployment scenarios, while simultaneously enhancing fairness by prioritizing cell-edge users who historically suffered from poor service quality. The theoretical benefits of NOMA have been validated through extensive research, leading to its standardization in 3GPP Release 15 and beyond under the name "Multi-User Superposition Transmission" (MUST), positioning it as a key enabling technology for 5G-Advanced and 6G networks. The fundamental principle lies in exploiting the channel gain disparity between users: weaker users receive higher transmit power to compensate for their poor channel conditions, while stronger users with better channel quality are allocated lower power but can still decode their signals reliably after performing SIC on the weaker users' signals.

\section{The Critical Challenge of User Pairing}

While NOMA's theoretical benefits are well-established and its advantages over OMA have been demonstrated through both analytical studies and system-level simulations, practical deployment faces a critical optimization problem that directly determines system performance: user pairing or clustering. Given $N$ users in a cell, the base station must intelligently decide which users should share each Physical Resource Block, a decision that profoundly impacts multiple aspects of system operation. Suboptimal pairings increase inter-user interference, significantly reducing achievable data rates and diminishing the spectral efficiency gains that NOMA promises. Successful interference cancellation requires sufficient channel gain disparity between paired users, typically $\geq$ 8 dB, making SIC feasibility a hard constraint that must be satisfied for reliable operation. Poor pairing decisions can also create severe fairness issues, potentially starving weak users of throughput or wasting valuable resources on strong users who could have been paired more effectively.

The computational complexity of optimal pairing presents perhaps the most significant practical challenge. The number of possible user pairings grows exponentially with user count—specifically, the number of perfect matchings on $N$ users is $(2N-1)!! = 1 \cdot 3 \cdot 5 \cdots (2N-1)$, which for $N=12$ users yields approximately $1.1 \times 10^{13}$ possible combinations. Evaluating each pairing at 1 microsecond per evaluation would require approximately 316 years of computation time, clearly impractical for real-time resource allocation in 5G systems operating with 1 millisecond Transmission Time Intervals. This combinatorial explosion necessitates intelligent algorithms that can find near-optimal solutions within stringent time constraints.

Traditional approaches to user pairing have evolved along several distinct paths, each offering different trade-offs between performance and computational complexity. Heuristic methods employ simple rules such as pairing the strongest user with the weakest user, offering computational efficiency with $O(N \log N)$ complexity but sacrificing optimality and adaptability to varying channel conditions. Optimization-based methods leverage graph matching algorithms such as the Blossom algorithm for maximum-weight general graph matching or the Hungarian algorithm for bipartite matching, achieving near-optimal performance but suffering from $O(N^3)$ computational complexity that becomes prohibitive for large-scale deployments. Reinforcement learning approaches attempt to learn adaptive policies through online interaction with the environment, but require extensive training episodes and may struggle with convergence in highly dynamic channel environments.

For large-scale deployments serving hundreds of users per cell, optimization methods become computationally prohibitive for meeting the real-time scheduling requirements of 5G New Radio (1 ms TTI), while simple heuristics sacrifice 20-30\% of potential throughput compared to optimal solutions. This performance-complexity trade-off represents a fundamental barrier to realizing NOMA's full potential in practical wireless systems.

\section{Graph Neural Networks: A Learning-Based Solution}

Recent advances in deep learning, particularly Graph Neural Networks (GNNs), offer a promising alternative paradigm that can potentially bridge the gap between heuristic efficiency and optimization-based performance. Graph Neural Networks can learn optimal pairing patterns from historical channel measurements and labeled training data without requiring hand-crafted rules or problem-specific heuristics. Once trained, they can generalize their learned policies to unseen scenarios and user distributions, adapting to variations in network topology, traffic patterns, and propagation environments. Most critically for practical deployment, GNN inference achieves $O(N \log N)$ or better computational complexity after offline training, enabling real-time operation even for large user populations.

GNNs naturally handle physical constraints such as SIC feasibility requirements and power budget limitations through careful architecture design and constraint-aware graph construction, avoiding the need for post-processing correction steps that can compromise solution quality. Among various GNN architectures, GraphSAGE (SAmple and aggreGatE), introduced by Hamilton et al. in 2017, enables inductive learning on graphs—predicting properties of previously unseen nodes and edges. This inductive capability is particularly valuable for wireless networks where user populations and network topology vary dynamically across time and location, allowing a single trained model to generalize across diverse deployment scenarios without requiring retraining.

\section{Literature Review: Evolution of NOMA Optimization}

The research literature on NOMA has evolved substantially over the past decade, progressing from theoretical analysis of two-user systems to sophisticated machine learning approaches for large-scale networks. Understanding this evolution provides essential context for the contributions of this thesis and highlights the research gaps that motivate our work.

\subsection{Foundations of Power-Domain NOMA}

The foundational concepts of power-domain NOMA were established through a series of influential works that demonstrated its superiority over orthogonal access schemes. Islam et al. (2017) provided a comprehensive survey of NOMA techniques, demonstrating through both analytical derivations and system-level simulations that NOMA can achieve 30-50\% capacity improvement over OMA in dense deployment scenarios. Their work established key insights: NOMA gains are maximized when users have significant channel gain disparity, fairness can be enhanced through careful power allocation favoring weak users, and SIC performance critically depends on the power allocation coefficient and received signal-to-interference-plus-noise ratio (SINR) levels. These fundamental principles guide NOMA system design and user pairing strategies.

Ding et al. (2014) provided rigorous theoretical analysis of NOMA's achievable rate region, deriving closed-form expressions for sum-rate under various power allocation strategies and establishing fundamental trade-offs between system throughput and user fairness. Their work demonstrated that NOMA's performance gains stem fundamentally from allowing controlled interference rather than eliminating it entirely, a paradigm shift from traditional interference-avoidance approaches. They also identified the critical role of channel state information quality in realizing these gains, motivating research into practical CSI acquisition and feedback mechanisms.

\subsection{Evolution of User Pairing Strategies}

Early NOMA research employed simple heuristic rules for user pairing, treating it as a secondary concern compared to power allocation optimization. Static pairing sorts users by channel gain and pairs the strongest with the weakest, maximizing channel disparity to facilitate SIC but ignoring spatial correlation and potentially creating suboptimal interference conditions. Balanced pairing divides users into upper and lower halves by channel quality and pairs one from each half, improving fairness metrics at a slight throughput cost but still relying on a predetermined rule that cannot adapt to channel dynamics.

Su et al. (2017) conducted systematic analysis of these heuristic methods, demonstrating that static pairing achieves 70-80\% of optimal throughput with $O(N \log N)$ complexity but exhibits significant performance degradation under non-uniform user distributions or correlated shadowing. Their results revealed that while heuristics provide computational efficiency, the gap to optimal performance grows substantially as network conditions deviate from idealized scenarios with independent, identically distributed channel gains.

Recognition that user pairing fundamentally determines NOMA performance motivated researchers to formulate it as a rigorous optimization problem. Hoang et al. (2019) proposed joint power allocation and user pairing via alternating optimization, solving the power allocation problem in closed form for fixed pairings and iterating over discrete pairing decisions. While convergence to a local optimum is guaranteed under mild conditions, computational cost scales as $O(N^2$ iterations$)$ and the solution quality depends heavily on initialization, limiting practical applicability for large networks.

\subsection{Graph-Theoretic Optimization Approaches}

Recognizing user pairing as fundamentally a matching problem on graphs, researchers adopted graph-theoretic formulations that enable polynomial-time algorithms with performance guarantees. Jia et al. (2018) formulated users as nodes in a general graph with edge weights representing achievable sum-rates for each potential pairing, then applied the Blossom algorithm developed by Edmonds (1965) to find the maximum-weight perfect matching. This approach achieves theoretical optimality for the sum-rate maximization objective but requires $O(N^3)$ computational complexity, limiting its scalability to large networks. Nevertheless, Blossom matching has become a standard benchmark against which other methods are compared, representing the gold standard for pairing quality.

To reduce computational complexity while maintaining near-optimal performance, several researchers proposed bipartite graph formulations where users are divided into two disjoint sets—typically "strong" and "weak" based on channel gain quantiles—creating a bipartite graph structure. The Hungarian algorithm, originally developed by Kuhn (1955) for the assignment problem, solves bipartite maximum-weight matching in $O(N^3)$ worst-case complexity but often achieves tighter practical bounds of $O(N^2 \log N)$ for sparse graphs. Zhang et al. (2023) extended this framework to incorporate proportional fairness weights, balancing throughput and equity by adjusting edge weights to reflect both sum-rate and individual user rate considerations. Their simulations demonstrated that fairness-aware bipartite matching achieves 90-92\% of optimal throughput while significantly improving minimum user rate compared to throughput-only optimization.

Köse et al. (2021) introduced an alternative graph-theoretic perspective by modeling user pairing through conflict graphs, where edges represent incompatible pairings that violate SIC constraints due to insufficient channel gain disparity or geometric conflicts. They proved that NOMA conflict graphs are claw-free under realistic channel conditions, enabling polynomial-time solutions to the Maximum Weighted Independent Set (MWIS) problem that otherwise would be NP-hard for general graphs. This elegant theoretical result demonstrates that NOMA's physical structure induces mathematical properties that can be exploited for efficient optimization.

Extending graph-based pairing to specialized scenarios, Mishra et al. (2021, 2022) developed bipartite matching algorithms for NB-IoT NOMA systems with unique constraints. Their Minimum-Weight Full Matching with Pruning (MWFMP) algorithm addresses connection density maximization subject to per-PRB power constraints, while the Connection Throughput Maximization with Bipartite Matching (CTMBM) scheme incorporates device deadlines and priority classes for grant-based uplink transmission. These works demonstrate that graph matching frameworks can be adapted to incorporate diverse system constraints and objectives beyond simple sum-rate maximization, achieving near-optimal performance (92-95\% of exhaustive search) with acceptable complexity for moderately sized networks ($N < 50$).

\subsection{Deep Learning for Wireless Resource Allocation}

The wireless communications community has increasingly embraced deep learning as a powerful tool for solving complex optimization problems that resist traditional analytical approaches. Sun et al. (2018) pioneered the "Learning to Optimize" paradigm for power control by training deep neural networks on optimal solutions generated by convex optimization solvers. Their key insight was that while solving convex optimization problems is computationally expensive, the mapping from problem instance (channel gains, QoS requirements) to optimal solution exhibits structure that neural networks can learn efficiently. Their trained networks achieved 95\% optimal performance with approximately 1000× speedup compared to iterative optimization, enabling real-time power control for spectrum sharing problems.

He et al. (2020) extended supervised learning approaches to multi-cell scenarios, demonstrating that learned policies exhibit remarkable generalization capabilities through transfer learning. Networks trained on specific cell geometries and user densities successfully generalized to significantly different configurations, suggesting that deep learning captures fundamental resource allocation principles rather than merely memorizing training scenarios. This generalization property is critical for practical deployment where network conditions vary continuously.

Reinforcement learning offers an alternative learning paradigm that does not require labeled training data from optimization solvers. Jiang et al. (2021) applied Deep Q-Networks (DQN) to joint NOMA pairing and power allocation, formulating the problem as a Markov Decision Process where the agent observes channel states and selects pairing actions to maximize cumulative discounted throughput. Their agent learned effective policies through trial-and-error interaction with a network simulator, achieving 88\% of optimal throughput after approximately 10,000 training episodes. However, reinforcement learning suffers from several practical limitations including sample inefficiency requiring extensive simulation time, sensitivity to reward function design where poor shaping can prevent convergence, and limited theoretical understanding of convergence guarantees and generalization properties.

Naparstek et al. (2019) demonstrated the benefits of multi-task learning (MTL) for wireless optimization, proposing a deep reinforcement learning framework that jointly optimizes power control and user association. By sharing learned representations across related tasks, MTL improves sample efficiency and generalization compared to training separate models for each task. For NOMA systems, MTL is particularly relevant since pairing decisions, power allocation, and rate prediction are inherently coupled—optimal pairings depend on achievable rates, which in turn depend on power allocation coefficients, creating circular dependencies that MTL can potentially resolve through joint optimization.

\subsection{Graph Neural Networks in Wireless Communications}

Graph Neural Networks represent a natural framework for wireless optimization problems that exhibit inherent graph structure such as interference relationships, network topology, and user proximity. The foundational GNN architectures include Graph Convolutional Networks (GCN) introduced by Kipf & Welling (2017) which apply spectral convolutions via normalized adjacency matrices, Graph Attention Networks (GAT) developed by Veličković et al. (2018) which incorporate attention mechanisms for adaptive neighbor weighting, and GraphSAGE proposed by Hamilton et al. (2017) which enables inductive learning through neighborhood sampling and aggregation. GraphSAGE's inductive capability is particularly valuable for wireless networks where user populations and topologies vary dynamically, allowing a single trained model to generalize to unseen network configurations without requiring retraining.

Recent applications demonstrate GNN effectiveness across diverse wireless optimization problems. Eisen et al. (2020) modeled interference networks as conflict graphs where nodes represent transmitters and edges indicate significant mutual interference, using Graph Convolutional Networks to learn distributed power allocation policies that achieve 95\% of centralized optimal solutions while operating in a fully distributed manner with each transmitter making local decisions based only on neighborhood information. Rusek et al. (2020) developed RouteNet, a GNN architecture that predicts network performance metrics including delay and jitter from routing configurations, enabling fast "what-if" analysis for network planning without requiring time-consuming discrete-event simulation.

Chowdhury et al. (2021) demonstrated an elegant connection between iterative optimization algorithms and GNN architectures by unfolding the Weighted Minimum Mean Square Error (WMMSE) algorithm for MIMO beamforming into GNN layers. Each iteration of WMMSE corresponds to one GNN layer, and the algorithm's message-passing structure maps naturally to graph neural network operations. By unfolding a fixed number of iterations and treating operation parameters as learnable weights, they created a differentiable architecture that combines optimization theory's principled structure with neural networks' representation learning capabilities, achieving superior performance to both purely optimization-based and purely learning-based approaches.

Shen et al. (2021) formulated wireless link scheduling as edge classification on interference graphs, where each edge represents a potential link activation and the goal is to select a subset that maximizes throughput while respecting interference constraints. Their GNN architecture processes the graph structure to predict which links should be activated, achieving near-optimal throughput with $O(E \log E)$ complexity compared to $O(E^2)$ for traditional greedy algorithms and exponential complexity for optimal exhaustive search.

\subsection{Graph Neural Networks for NOMA: Current State}

Despite the promising applications of GNNs to various wireless optimization problems, their application to NOMA user pairing remains relatively unexplored with only limited prior work. Zhao et al. (2020) proposed a GNN-based approach for NOMA resource allocation in vehicular networks, treating vehicles as graph nodes and potential communication links as edges. Their model jointly learned user clustering and power allocation, achieving 88\% of optimal performance with approximately 50× speedup compared to exhaustive search. However, this work focused specifically on vehicle-to-vehicle communication scenarios with mobility-specific features such as velocity, trajectory prediction, and Doppler effects, and did not explicitly incorporate physical constraints like SIC feasibility into the graph construction process, instead relying on post-processing to filter invalid pairings.

Wang et al. (2021) explored Graph Attention Networks for uplink NOMA in massive IoT scenarios where thousands of devices compete for access. Their GAT architecture learned to prioritize critical devices requiring low-latency transmission while deferring delay-tolerant traffic. While achieving good throughput and fairness metrics, the approach required labeled training data generated through exhaustive search over all possible device schedules, severely limiting scalability to large device populations where exhaustive search becomes computationally infeasible even for offline data generation.

\subsection{Identified Research Gaps}

Based on comprehensive analysis of the existing literature, several critical gaps emerge that motivate the research presented in this thesis. First, there exists a fundamental scalability-optimality trade-off where no existing method achieves both near-optimal performance (>95\% of theoretical optimum) and low computational complexity ($O(N \log N)$ or better). Graph matching algorithms provide optimality guarantees but require $O(N^3)$ complexity, while heuristics achieve linear complexity but sacrifice 20-30\% throughput. Second, most deep learning approaches ignore physical constraints such as SIC feasibility requirements, channel gain disparity thresholds, and angular separation requirements, instead training unconstrained models and applying post-hoc filtering that can significantly degrade solution quality and reduce the fraction of users that can be successfully paired.

Third, existing methods typically decouple user pairing from power allocation, optimizing them in separate sequential stages rather than jointly optimizing the inherently coupled decisions. This decomposition provably yields suboptimal solutions since the optimal pairing depends on achievable rates, which depend on power allocation coefficients, creating circular dependencies that sequential optimization cannot resolve. Fourth, many works employ simplified channel models such as distance-dependent path loss with fixed exponents rather than using standardized 3GPP models incorporating realistic effects including line-of-sight probability, log-normal shadowing, frequency-dependent path loss, and environment-specific parameters. This limits the relevance of results for practical deployment. Fifth, limited empirical analysis exists regarding generalization performance—the ability of learned models to maintain performance on scenarios substantially different from training data, which is critical for practical deployment where network conditions vary continuously.

\section{Problem Formulation and Thesis Objectives}

Having established the importance of user pairing in NOMA systems and identified critical gaps in existing approaches, we now formally state the optimization problem and outline the specific objectives and contributions of this thesis.

\subsection{System Description and Assumptions}

We consider a single-cell downlink PD-NOMA system serving $N$ users denoted by the set $\mathcal{N} = \{1, 2, \ldots, N\}$. The base station is equipped with a single transmit antenna operating at carrier frequency $f_c = 3.5$ GHz (5G NR band n78) with total transmit power $P = 1.0$ W (30 dBm). Each user equipment is equipped with a single receive antenna and has perfect SIC capability, meaning it can perfectly decode and cancel interfering signals provided the SINR exceeds a minimum threshold $\gamma_{\text{th}}$. We assume perfect Channel State Information (CSI) is available at the base station through uplink feedback, and channels exhibit quasi-static flat fading that remains constant over each resource allocation interval (1 ms TTI) but varies independently across intervals.

The channel coefficient between the base station and user $i$ is modeled according to the 3GPP TR 38.901 Urban Macro (UMa) specification as $h_i = \sqrt{PL_i \cdot SF_i} \cdot g_i$, where $PL_i$ represents distance-dependent path loss including both line-of-sight and non-line-of-sight components, $SF_i$ represents log-normal shadow fading with standard deviation 4-6 dB depending on LOS condition, and $g_i \sim \mathcal{CN}(0,1)$ represents small-scale Rayleigh fading. The channel power gain is $|h_i|^2 = PL_i \cdot SF_i \cdot |g_i|^2$, and we order users such that $|h_1|^2 \leq |h_2|^2 \leq \cdots \leq |h_N|^2$ from weakest to strongest channel gain.

\subsection{NOMA Transmission and Rate Analysis}

For a two-user cluster pairing user $i$ (weak) with user $j$ (strong) where $|h_i|^2 < |h_j|^2$, the base station transmits the superposed signal $s = \sqrt{\alpha P} \cdot x_i + \sqrt{(1-\alpha) P} \cdot x_j$ on a single Physical Resource Block with bandwidth $B = 20$ MHz, where $x_i$ and $x_j$ are unit-power information symbols and $\alpha \in (0,1)$ is the power allocation coefficient. Following NOMA principles, we allocate higher power to the weaker user, thus $\alpha > 0.5$ typically.

The weaker user $i$ directly decodes its signal treating user $j$'s signal as interference, achieving SINR of $\text{SINR}_i = \frac{\alpha P |h_i|^2}{(1-\alpha) P |h_i|^2 + \sigma^2}$ where $\sigma^2 = kTB$ is the thermal noise power with $k$ being Boltzmann's constant and $T$ the noise temperature. The corresponding achievable rate using Shannon capacity formula is $R_i = B \log_2(1 + \text{SINR}_i)$ in bits per second.

The stronger user $j$ employs successive interference cancellation, first decoding user $i$'s signal (treating its own signal as interference) with SINR of $\text{SINR}_{j \to i} = \frac{\alpha P |h_j|^2}{(1-\alpha) P |h_j|^2 + \sigma^2}$. After successfully decoding $x_i$, user $j$ performs SIC to cancel user $i$'s contribution from the received signal, then decodes its own signal interference-free with SINR of $\text{SINR}_j = \frac{(1-\alpha) P |h_j|^2}{\sigma^2}$. The achievable rate for user $j$ is $R_j = B \log_2(1 + \text{SINR}_j)$.

For successful SIC operation, user $j$ must decode user $i$'s signal with sufficient reliability, requiring $\text{SINR}_{j \to i} \geq \gamma_{\text{th}}$ where $\gamma_{\text{th}}$ is the SIC decoding threshold typically set to 8 dB (6.31 in linear scale) based on practical receiver implementations. Substituting the SINR expression and manipulating algebraically, this constraint implies a minimum channel gain disparity requirement: $\frac{|h_j|^2}{|h_i|^2} \geq \gamma_{\text{th}}$ approximately for high SNR scenarios where $P|h_j|^2 \gg \sigma^2$. Thus, successful NOMA operation requires pairing users with sufficient channel quality difference, typically 8 dB or more.

\subsection{Optimization Problem Statement}

The joint user pairing and power allocation problem aims to maximize system sum-rate while ensuring SIC feasibility and satisfying matching constraints. Formally, let binary variables $x_{ij} \in \{0,1\}$ indicate whether users $i$ and $j$ are paired ($x_{ij} = 1$) or not ($x_{ij} = 0$), and let $\alpha_{ij} \in (0,1)$ denote the power allocation coefficient for pair $(i,j)$ if paired. The optimization problem is:

\begin{equation}
\begin{aligned}
\max_{x, \alpha} \quad & \sum_{i=1}^{N} \sum_{j=i+1}^{N} x_{ij} \cdot R_{\text{sum}}(i, j, \alpha_{ij}) \\
\text{subject to} \quad & \sum_{j=1, j \neq i}^{N} x_{ij} \leq 1, \quad \forall i \in \mathcal{N} \quad \text{(matching constraint)} \\
& x_{ij} \in \{0, 1\}, \quad \forall i, j \in \mathcal{N}, i \neq j \\
& \frac{|h_j|^2}{|h_i|^2} \geq \gamma_{\text{th}}, \quad \text{if } x_{ij} = 1, |h_i|^2 < |h_j|^2 \quad \text{(SIC constraint)} \\
& |\theta_i - \theta_j| \geq \Delta\theta_{\min}, \quad \text{if } x_{ij} = 1 \quad \text{(angular separation)} \\
& \alpha_{ij} \in (0, 1), \quad \forall i, j \in \mathcal{N}
\end{aligned}
\end{equation}

where $R_{\text{sum}}(i,j,\alpha_{ij}) = R_i + R_j$ is the sum-rate for pair $(i,j)$ with power allocation $\alpha_{ij}$, $\theta_i$ denotes the angular position of user $i$ in the cell, and $\Delta\theta_{\min}$ is the minimum required angular separation (typically 30 degrees) to ensure spatial decorrelation.

This optimization problem is a Mixed-Integer Nonlinear Program (MINLP) that is NP-hard in general. The binary pairing variables $x_{ij}$ create $\binom{N}{2}$ discrete decisions, while the continuous power allocation variables $\alpha_{ij}$ couple nonlinearly with the discrete decisions through the rate expressions. The number of possible matching configurations (including leaving users unpaired) grows as the double factorial $(2N-1)!! = 1 \cdot 3 \cdot 5 \cdots (2N-1)$, which for $N=12$ users yields approximately $1.1 \times 10^{13}$ possible matchings. Evaluating each matching at 1 microsecond per evaluation would require approximately 316 years of computation time, clearly demonstrating the impracticality of exhaustive search for even moderate user populations.

\subsection{Thesis Objectives and Contributions}

The primary objective of this thesis is to develop a scalable, near-optimal solution to the NOMA user pairing problem that achieves the performance of graph-theoretic optimization methods while attaining the computational efficiency of heuristic approaches. Specifically, we aim to design a Graph Neural Network framework that learns optimal pairing strategies from training data generated by the Blossom maximum-weight matching algorithm, then performs real-time inference with $O(N \log N)$ complexity while maintaining at least 95\% of optimal throughput. To accomplish this overarching goal, the thesis pursues the following specific objectives:

First, we implement realistic channel modeling using the 3GPP TR 38.901 Urban Macro (UMa) specification to generate high-fidelity Channel State Information datasets incorporating distance-dependent path loss with separate models for line-of-sight and non-line-of-sight propagation, log-normal shadow fading with appropriate standard deviations, Rayleigh small-scale fading for rich scattering environments, and realistic user spatial distributions within a circular cell. This ensures that all evaluation results reflect practical deployment scenarios rather than idealized theoretical models.

Second, we implement and comprehensively benchmark four representative traditional user clustering strategies to establish performance baselines: Static Pairing which pairs strongest users with weakest users to maximize channel disparity; Balanced Pairing which divides users into upper and lower halves by channel quality and pairs one from each half; Blossom Maximum-Weight Matching which applies Edmonds' algorithm to find the globally optimal pairing for sum-rate maximization; and Bipartite Graph Matching which formulates strong-weak pairing as a bipartite assignment problem solved via the Hungarian algorithm with proportional fairness weights. These baselines span the spectrum from simple heuristics to optimal algorithms, providing context for evaluating our proposed learning-based approach.

Third, we design and implement NOMA-GNN, a novel GraphSAGE-based deep learning architecture specifically tailored for NOMA user pairing. The framework treats users as nodes in a graph with CSI-derived features including channel power gain, distance from base station, and angular position, while potential user pairings correspond to edges whose presence is determined by SIC feasibility and angular separation constraints. The GNN learns message-passing representations through neighborhood aggregation that capture complex user relationships and channel dependencies, predicts pairwise compatibility scores indicating which user pairs would achieve high sum-rate if paired together, and employs greedy matching on predicted scores to construct valid pairing solutions with $O(N \log N)$ complexity. Crucially, we enforce physical constraints during graph construction rather than as post-processing, ensuring all proposed pairings are feasible.

Fourth, we develop a complete end-to-end learning pipeline encompassing data preparation with automated graph construction from raw CSV channel measurements, model training using optimal pairing labels from Blossom algorithm as supervision with carefully designed loss functions and regularization, real-time inference incorporating SIC verification and constraint checking, and comprehensive evaluation measuring multiple performance dimensions. The pipeline is designed for practical deployment with a lightweight model architecture containing approximately 128K parameters, suitable for implementation on edge computing platforms at base stations.

Fifth, we conduct extensive experimental evaluation comparing traditional and learning-based methods across multiple performance metrics including system throughput measured as aggregate sum-rate across all paired users, fairness quantified using Jain's fairness index and minimum user rate, computational complexity measured as wall-clock inference time and theoretical complexity scaling, pairing accuracy measured as agreement with optimal Blossom solutions, and generalization performance evaluated on scenarios with different user densities and channel statistics than the training set.

The main contributions of this thesis are threefold. First, we provide the first systematic comparison of traditional graph-based NOMA pairing methods using standardized 3GPP UMa channel models with realistic path loss, shadowing, and fading statistics, establishing credible benchmarks for future research. Second, we propose a novel Graph Neural Network architecture NOMA-GNN featuring constraint-aware graph construction that enforces SIC and angular separation requirements, multi-task learning jointly predicting pairing compatibility and achievable rates, inductive learning capability enabling generalization to unseen network configurations, and lightweight design with 128K parameters suitable for edge deployment. Third, we demonstrate superior performance-complexity trade-off achieving 96.5\% of optimal Blossom throughput (70.74 Gbps versus 73.30 Gbps) with 52.9× computational speedup (846 ms versus 44,749 ms), establishing practical viability for real-time deployment in dense networks serving 500+ users per cell.

\subsection{Scope and Limitations}

This research focuses specifically on downlink Power-Domain NOMA in single-cell Urban Macro scenarios representative of dense urban 5G deployments. We restrict attention to two-user pairing (cluster size of 2) as this configuration achieves the best balance between spectral efficiency gains and implementation complexity based on prior research. We assume perfect CSI is available at the base station through error-free uplink feedback, abstracting away channel estimation and feedback quantization issues that would be present in practical systems. The channel model assumes quasi-static fading where channel coefficients remain constant during each 1 ms resource allocation interval but vary independently across intervals, capturing block-fading behavior while simplifying the optimization problem by eliminating predictive elements.

Several important aspects lie beyond the scope of this thesis and represent opportunities for future work. We do not address multi-cell coordination and inter-cell interference management, which would be essential for deployment in cellular networks where adjacent cells create interference. Larger cluster sizes with more than 2 users per PRB are not considered, though they could potentially offer additional spectral efficiency at the cost of increased SIC complexity and error propagation. Imperfect CSI due to channel estimation errors, feedback quantization, and outdated measurements represents a significant practical concern not addressed in our idealized perfect CSI assumption. Dynamic channel prediction and user mobility modeling would be necessary for extending the framework to mobile scenarios where users move during the resource allocation interval. Uplink NOMA and full-duplex communication scenarios are not considered, focusing instead on the more widely studied downlink transmission. Finally, hardware implementation and over-the-air experimental validation would be required to confirm that the performance improvements demonstrated in simulation translate to real-world deployments with practical receiver implementations and RF impairments.

\subsection{Thesis Organization}

The remainder of this thesis is organized to provide a comprehensive treatment of the NOMA user pairing problem and our proposed solution. Chapter 2 describes the detailed system model including base station configuration, user equipment capabilities, the 3GPP TR 38.901 Urban Macro channel model implementation with path loss, shadowing and fading components, and the mathematical framework for NOMA signal transmission and rate analysis. Chapter 3 presents traditional user clustering methodologies including static pairing, balanced pairing, and graph-theoretic optimization approaches using Blossom and bipartite matching algorithms, along with computational complexity analysis for each method. Chapter 4 introduces the proposed NOMA-GNN framework detailing the graph construction procedure, GraphSAGE architecture design, training methodology including loss functions and optimization, and the inference pipeline for real-time pairing prediction.

Chapter 5 presents comprehensive results including implementation details and software architecture, experimental setup specifying datasets and hyperparameters, performance comparison between traditional and learning-based methods across throughput, fairness, and complexity metrics, ablation studies analyzing the contribution of individual components, and deployment analysis examining practical considerations for real-world implementation. Chapter 6 concludes the thesis with a summary of key findings, discussion of practical implications for NOMA deployment in 5G and beyond systems, and identification of promising directions for future research. Appendices provide complete source code listings to enable reproducibility, detailed dataset specifications describing the channel model parameters and data generation procedure, and hyperparameter configurations for the GNN model including architecture details and training settings.

This comprehensive structure ensures that readers can understand both the theoretical foundations and practical implementation of our proposed approach, evaluate its performance relative to existing methods, and extend the work in future research directions.
